% !TEX root = ../main.tex
\section{Introduction}

Neutrino oscillations are one of the clearest experimental indications that the Standard Model of particle physics must be extended, as established by the atmospheric-neutrino measurements of Super-Kamiokande and the solar-neutrino measurements of the Sudbury Neutrino Observatory \parencite{SuperKamiokande1998,SNO2002}. A neutrino produced with a definite flavour, such as an electron, muon or tau neutrino, propagates as a coherent superposition of mass eigenstates. The interference among these components gives rise to flavour transitions whose probabilities depend on the neutrino energy, the travelled distance, the mixing parameters and, when propagation occurs in matter, the electron density of the medium \parencite{Giunti2007,Wolfenstein1978,MikheyevSmirnov1985}.

Modern neutrino analyses require evaluating these probabilities many times. Oscillation probabilities appear inside flux predictions, detector simulations, likelihood scans and global fits. For this reason, the numerical implementation of the oscillation formalism is not a secondary detail: it directly affects the feasibility of high-statistics studies and of exploratory scans over large parameter spaces. Existing tools such as PEANUTS \parencite{Gonzalo2024} provide compact routines for neutrino propagation, matter mixing and solar-neutrino probabilities. However, array-based implementations are not always optimal for current workflows involving GPUs, TPUs, automatic differentiation and batched execution.

This thesis studies and documents the development of \texttt{tpeanuts}, a PyTorch \parencite{Paszke2019}-based reproduction and extension of the PEANUTS neutrino-oscillation toolkit. The starting point is a validated, PyTorch-native translation of the core of PEANUTS, three-flavour vacuum and matter propagation and solar-neutrino probabilities, that preserves its physical conventions while expressing every operation as a tensor computation. This makes it possible to evaluate many oscillation configurations in parallel, to move computations between CPU and GPU with minimal changes, and to connect the resulting probabilities to gradient-based parameter estimation.

Built on top of this validated core, \texttt{tpeanuts} extends substantially PEANUTS: a modular architecture adds atmospheric-neutrino propagation through the Earth, several interchangeable solar and Earth matter/density models, beyond-the-Standard-Model scenarios (non-standard interactions and sterile neutrinos), and a full detector-level forward-folding and statistical-inference layer, validated against real, published data from several experiments. These extensions, together with the underlying oscillation engine, are released as a public, open-source Python package, with installation and usage documentation (\cref{sec:supplementary-material}), so that other researchers can reuse and extend it directly rather than reimplementing it from the physics described here.

\subsection{Motivation}

The main motivation for \texttt{tpeanuts} is the need for fast, repeated neutrino-oscillation evaluation, both for large-scale simulation (atmospheric-flux studies, detector-response studies, BSM scans) and, ultimately, for fitting real data: a gradient-based fit to the published spectrum of a detector requires the oscillation probability, the detector response and the likelihood to be evaluated, and differentiated, at every optimiser step. PyTorch provides a natural environment for this because it combines vectorised tensor operations, GPU support and automatic differentiation, letting a single implementation serve fast forward simulation and gradient-based inference alike.

The original PEANUTS implementation is tied to a specific solar model and a fixed treatment of matter effects, which limits its use where the sensitivity of a result to these choices must be tested. Building \texttt{tpeanuts} on a flexible tensor framework makes it possible to add such options, including alternative solar and Earth density models, perturbative treatments of the matter potential, and new-physics scenarios such as non-standard interactions (NSI) and sterile neutrinos, in a modular way, without compromising the validated core of the oscillation formalism. This same modularity is what makes it practical to add a detector-response layer (target composition, energy resolution, backgrounds, efficiencies) and an inference layer (likelihood construction, gradient-based fitting, uncertainty estimation) on top of the oscillation engine, and to reuse both across experiments as different as a reactor (Daya Bay), a liquid-scintillator solar detector (Borexino), a large-volume atmospheric detector (IceCube DeepCore) and a heavy-water Cherenkov detector (SNO).

Beyond the physics case for building it, three complementary features make \texttt{tpeanuts} useful to the wider neutrino-phenomenology community rather than only to this thesis. It is fast and scalable for large batched calculations: because the propagation engine operates on tensors rather than on scalars, thousands of oscillation configurations, spanning energies, baselines or points in a BSM parameter space, are evaluated in a single call instead of one at a time, with the advantage growing with batch size. It provides a common, flexible framework for solar, atmospheric and BSM neutrino studies that would otherwise require separate, purpose-built codes. And, being built on PyTorch, it is differentiable end to end, so oscillation probabilities can be connected directly to gradient-based inference and machine-learning applications rather than only to finite-difference or sampling-based methods. Combined with the fact that the code is public, on GitHub, with installation instructions and a per-module companion documentation set, this positions as a reusable scientific tool rather than a single-purpose thesis codebase.

Delivering on this motivation is only credible if every result is validated, so this thesis validates \texttt{tpeanuts} along three complementary fronts. Internal validation checks the physics of the implemented neutrino-oscillation phenomena, verifying that independent numerical strategies for the same propagation problem agree with each other within the expected approximation. External validation compares \texttt{tpeanuts} against the original PEANUTS implementation and against \texttt{nuSQuIDS} \parencite{nuSQuIDS}, an independent, widely used neutrino-propagation toolkit, at the level of both intermediate quantities and final probabilities. Finally, established experimental and observational results from the literature, from the solar day--night asymmetry to real detector data, are reproduced directly. Combining these three fronts is what makes it possible to obtain a tool that is simultaneously fast, reliable and versatile.

\subsection{Objectives}

The general objective of this Master's Thesis is to develop and document \texttt{tpeanuts} as a modular, extensible and publicly available PyTorch framework for neutrino-oscillation phenomenology, built on a validated translation of the core of PEANUTS and extended into a detector-level statistical-inference framework validated against real experimental data. The specific objectives are:

\begin{itemize}
  \item to review the three-flavour oscillation formalism used by PEANUTS, including vacuum propagation and matter effects, and to translate its numerical core to PyTorch tensors while preserving physical conventions, units and sign choices;
  \item to design a modular architecture in which vacuum/solar/Earth/atmosphere models, matter and perturbative treatments, and new-physics scenarios (NSI, sterile neutrinos) can be selected and reused consistently across the application;
  \item to define a validation strategy comparing \texttt{tpeanuts} against PEANUTS and \texttt{nuSQuIDS} at the level of intermediate quantities and final probabilities, complemented by the reproduction of well-known theoretical and experimental results and by a CPU/GPU benchmarking methodology;
  \item to implement a detector-agnostic forward-folding layer (cross sections, target composition, energy response, efficiency, backgrounds) and to specialise it for real experiments;
  \item to implement a gradient-based statistical-inference layer (Poisson and Gaussian likelihoods, LBFGS fitting, Laplace/Fisher uncertainty estimation, likelihood scans) built on PyTorch autograd;
  \item to validate the full detector-inference chain against real, published data, with Daya Bay as the primary case study, and to document its application to Borexino, IceCube DeepCore and SNO;
  \item to release \texttt{tpeanuts} publicly, as a documented, installable, open-source package intended for reuse by the neutrino-phenomenology community, rather than as a closed thesis-specific codebase.
\end{itemize}

\subsection{Document structure}

Section~2 introduces the theoretical framework of neutrino masses, mixing, oscillations, matter effects and the statistical-inference formalism. Section~3 describes the methodology used to translate and validate the numerical core, extend it into a modular propagation framework and into beyond-the-Standard-Model scenarios, validate the resulting inference chain against real data, and release the code publicly. Section~4 presents the technical development of \texttt{tpeanuts}, including its detector and inference layers. Section~5 first establishes the numerical reliability and computational usefulness of the implementation through validation and benchmarking against PEANUTS and \texttt{nuSQuIDS}, then demonstrates capabilities that go beyond a translation of PEANUTS, atmospheric-neutrino propagation through the Earth and flexible matter/density models, followed by the NSI and sterile-neutrino extensions, then the phenomenological application of the full framework to real experimental data, headlined by the Daya Bay near/far detector-level case study, and closes with the public repository underlying every result. Finally, Section~6 summarises the conclusions and outlines future work.
